\chapter{Conclusion}
\label{ch_conclusion}


% \textbf{Overall Motivation and goal.} 
% ----------------------------------------------------------------
The Class I disk IRS 63 was studied using multi-wavelength polarization observations from ALMA at \bfourNS, \bfiveNS, \bsixNS, and \bsevenNS. The main goal of this project was to constrain the dust grain size in the disk to look for evidence of dust grain growth during this early evolutionary stage. The observed polarization morphology and dust scattering models provided constraints on the dust grain properties, giving insight into the early stages of planet formation. 
% ----------------------------------------------------------------


% ----------------------------------------------------------------
% \textbf{Polarization Morphology change with Wavelength and spatially} 
% ----------------------------------------------------------------
The observed polarization morphology shows a clear wavelength dependence, with patterns at longer wavelengths following an azimuthal pattern, shorter wavelengths being uniformly aligned with the disk minor axis, and intermediate patterns at intermediate wavelengths. These patterns suggest that dust self-scattering dominates at shorter wavelengths, while another mechanism may dominate at longer wavelengths (e.g. radiative grain alignment). This wavelength dependence was also seen in HL~Tau in previous studies. The polarization morphology also varies spatially across the disk, with the azimuthal pattern dominating in the outer region of the disk, while the uniform pattern dominates closer to the centre. Based on the expectation that dust self-scattering is more efficient at higher optical depth, this spatial dependence could be due to scattering being more efficient at the optically thick centre of the disk, and less efficient at the optically thin outer edges. The observational analysis also identified a streamer and a potential background galaxy in the observations. These findings did not affect the main analysis or conclusions of the project.


Dust self-scattering models were used to determine if the observed polarization fractions from the central region could be explained by dust self-scattering, and were used to constrain the maximum dust grain size. Scattering calculations for spherical grains were performed using \texttt{DSHARP} dust models for various porosities, and calculations for compact irregular grains utilized \texttt{glitterin}. Models for both grain shapes have similar assumptions, including the DSHARP dust composition and the same grain size distribution, to investigate how grain shape, size and porosity affect polarization from scattering and constrain these properties.  
% ----------------------------------------------------------------





% ----------------------------------------------------------------
% \textbf{Evidence for dust grain growth. } 
% ----------------------------------------------------------------
The different models constrain the maximum grain size to be \sbest for spherical grains with a moderate porosity ($f$ = 0.5) and \ibest for irregular grains. These constrained dust grain sizes in IRS 63 are substantially larger than those expected in the ISM or the natal dense core by at least two orders of magnitude. The constrained grain sizes are also relatively consistent with grain sizes found in other studies of disks, suggesting that grain sizes of a few hundred microns are common in disks, even at earlier stages.
% ----------------------------------------------------------------






% ----------------------------------------------------------------
% \textbf{Final statement:}
% ----------------------------------------------------------------
Overall, this study provides evidence suggesting that dust grain growth, an initial stage in planet formation, may have occurred in IRS 63 by the Class I stage. This has applications for planet formation as a whole, as planet formation may be occurring earlier than currently included in some models. To be more certain and explore dependencies, this project would be best extended by applying the methods to other Class I disks. Additionally, the study could be extended to other regions in the disk. Future improvements to ALMA will potentially make a polarization population study of disks possible, providing a great opportunity to further investigate dust grain growth and improve our understanding of planet formation.
% ----------------------------------------------------------------
